\documentclass[12pt]{article}
\usepackage{amsmath, amssymb}
\usepackage{geometry}
\usepackage{enumitem}
\geometry{margin=1in}

\title{CSE 400: Fundamentals of Probability in Computing}
\author{Rudra Mehta \\ AU2440268}
\date{26/2/26}

\begin{document}

\maketitle

\begin{center}
\textbf{Lecture 14 – Tutorial 2 Solutions} \\
(SEAS, Ahmedabad University)
\end{center}

\vspace{0.5cm}

This scribe reconstructs all solutions exactly in the order and structure presented in the tutorial document. All definitions, assumptions, derivations, and intermediate steps are written explicitly for exam revision.

\hrule
\vspace{0.5cm}

\section*{Q1}

\subsection*{Problem}

A point is chosen at random on a line segment of length $L$. Interpret this statement and find the probability that the ratio of the shorter to the longer segment is less than $\frac{1}{4}$.

\subsection*{Step 1: Mathematical Interpretation}

Let the point be chosen at a distance $x$ from one end.

\[
0 < x < L
\]

The two segments formed are:

\begin{itemize}
\item Length $x$
\item Length $L - x$
\end{itemize}

The ratio of the shorter to the longer segment is:

\[
R = \frac{\min(x, L-x)}{\max(x, L-x)}
\]

We must compute:

\[
P\left(R < \frac{1}{4}\right)
\]

\subsection*{Step 2: Case Analysis Based on Midpoint}

We divide logically at the midpoint $x = \frac{L}{2}$.

\subsubsection*{Case 1: $x \le \frac{L}{2}$}

Shorter segment = $x$

Longer segment = $L-x$

\[
\frac{x}{L-x} < \frac{1}{4}
\]

Multiply both sides:

\[
4x < L - x
\]

\[
5x < L
\]

\[
x < \frac{L}{5}
\]

\subsubsection*{Case 2: $x > \frac{L}{2}$}

Shorter segment = $L - x$

Longer segment = $x$

\[
\frac{L-x}{x} < \frac{1}{4}
\]

Multiply both sides:

\[
4L - 4x < x
\]

\[
5x > 4L
\]

\[
x > \frac{4L}{5}
\]

\subsection*{Step 3: Valid Region}

\[
x \in \left(0, \frac{L}{5}\right) \cup \left(\frac{4L}{5}, L\right)
\]

Total valid length:

\[
\frac{L}{5} - 0 + \left(L - \frac{4L}{5}\right)
= \frac{L}{5} + \frac{L}{5}
= \frac{2L}{5}
\]

\subsection*{Step 4: Probability}

Since the point is uniformly chosen over length $L$:

\[
P = \frac{\text{valid region}}{\text{total region}}
= \frac{2L/5}{L}
= \frac{2}{5}
\]

\hrule
\vspace{0.5cm}

\section*{Q2}

\subsection*{Given}

White region (W):
\[
X|W \sim N(4,4), \quad \sigma_1 = 2
\]

Black region (B):
\[
X|B \sim N(6,9), \quad \sigma_2 = 3
\]

\[
P(B) = \alpha, \quad P(W) = 1-\alpha
\]

Observation: $X = 5$

We require equal probability of error regardless of conclusion:

\[
P(B|X=5) = P(W|X=5) = \frac{1}{2}
\]

\subsection*{Step 1: Apply Bayes’ Theorem}

\[
P(B|X=5)
=
\frac{\alpha f(5|B)}
{\alpha f(5|B) + (1-\alpha) f(5|W)}
= \frac{1}{2}
\]

\subsection*{Step 2: Normal Density}

\[
f(x) = \frac{1}{\sigma\sqrt{2\pi}}
e^{-\frac{(x-\mu)^2}{2\sigma^2}}
\]

Compute densities:

\[
f(5|B) = \frac{1}{3\sqrt{2\pi}} e^{-1/18}
\]

\[
f(5|W) = \frac{1}{2\sqrt{2\pi}} e^{-1/8}
\]

\subsection*{Step 3: Substitute}

After substitution and simplification:

\[
\frac{3\alpha}{3\alpha + 2(1-\alpha)e^{-5/72}}
= \frac{1}{2}
\]

\subsection*{Step 4: Solve}

\[
6\alpha = 3\alpha + 2(1-\alpha)e^{-5/72}
\]

Using $e^{-5/72} \approx 0.93$:

\[
3\alpha = 1.86 - 1.86\alpha
\]

\[
4.86\alpha = 1.86
\]

\[
\alpha \approx 0.3827
\]

\[
\boxed{\alpha \approx 0.38}
\]

\hrule
\vspace{0.5cm}

\section*{Q3}

\subsection*{(a) Uniform Case}

\[
X \sim \text{Uniform}(0,A)
\]

\[
f_X(x) = \frac{1}{A}
\]

We minimize:

\[
E[|X-a|]
=
\int_0^a (a-x)\frac{1}{A}dx
+
\int_a^A (x-a)\frac{1}{A}dx
\]

After evaluation and differentiation:

\[
\frac{d}{da}
\left(
\frac{2a^2 - 2aA + A^2}{2A}
\right)
=
\frac{1}{2A}(4a-2A)
\]

Set equal to zero:

\[
4a = 2A
\]

\[
a = \frac{A}{2}
\]

Station at midpoint.

\subsection*{(b) Exponential Case}

\[
X \sim \text{Exponential}(\lambda)
\]

\[
f_X(x) = \lambda e^{-\lambda x}
\]

\[
E[|X-a|]
=
\int_0^a (a-x)\lambda e^{-\lambda x} dx
+
\int_a^\infty (x-a)\lambda e^{-\lambda x} dx
\]

Differentiate:

\[
\frac{d}{da}E[|X-a|]
=
1 - 2e^{-\lambda a}
\]

Set equal to zero:

\[
2e^{-\lambda a} = 1
\]

\[
e^{-\lambda a} = \frac{1}{2}
\]

\[
-\lambda a = -\ln 2
\]

\[
a = \frac{\ln 2}{\lambda}
\]

\hrule
\vspace{0.5cm}

\section*{Q4}

Let mixture of exponentials.

\[
P(X>t) = p_1 e^{-\lambda_1 t} + p_2 e^{-\lambda_2 t}
\]

\[
P(X>t+s)
=
p_1 e^{-\lambda_1(t+s)}
+
p_2 e^{-\lambda_2(t+s)}
\]

Conditional probability:

\[
P(X>t+s|X>t)
=
\frac{
p_1 e^{-\lambda_1(t+s)} + p_2 e^{-\lambda_2(t+s)}
}{
p_1 e^{-\lambda_1 t} + p_2 e^{-\lambda_2 t}
}
\]

\hrule
\vspace{0.5cm}

\section*{Q5}

\[
X \sim \text{Poisson}(\lambda_1)
\]
\[
Y \sim \text{Poisson}(\lambda_2)
\]

We compute:

\[
P(X=k|X+Y=n)
=
\frac{P(X=k)P(Y=n-k)}{P(X+Y=n)}
\]

After substitution and simplification:

\[
\binom{n}{k}
\left(\frac{\lambda_1}{\lambda_1+\lambda_2}\right)^k
\left(\frac{\lambda_2}{\lambda_1+\lambda_2}\right)^{n-k}
\]

\hrule
\vspace{0.5cm}

\section*{Q6}

Given:

\[
X \sim \text{Uniform}[-2,2]
\]

Transformation:

\[
g(x)=
\begin{cases}
x, & x\in[-2,-1] \\
0, & x\in(-1,1) \\
x, & x\in[1,2]
\end{cases}
\]

\subsection*{CDF (Piecewise)}

\begin{itemize}
\item $y < -2$: \quad $0$
\item $-2 \le y < -1$: \quad $\dfrac{y+2}{4}$
\item $-1 \le y < 0$: \quad $\dfrac{1}{4}$
\item At $y=0$: point mass
\end{itemize}

\[
P(Y=0)=\frac{2}{4}=\frac{1}{2}
\]

\[
F_Y(0)=\frac{3}{4}
\]

\begin{itemize}
\item $0<y<1$: \quad $\dfrac{3}{4}$
\item $1\le y\le2$: \quad $\dfrac{y+2}{4}$
\item $y>2$: \quad $1$
\end{itemize}

\subsection*{PDF}

\[
f_Y(y)=\frac{1}{4}, \quad y\in(-2,-1)\cup(1,2)
\]

Point mass:

\[
P(Y=0)=\frac{1}{2}
\]

\hrule
\vspace{0.5cm}

\section*{Q7}

Let number of ups be $k$.

\[
S_{1000}=S_0 u^k d^{1000-k}
\]

Condition:

\[
u^k d^{1000-k} \ge 1.30
\]

Take logs:

\[
k(\ln u - \ln d) + 1000\ln d \ge \ln(1.30)
\]

Using numerical values:

\[
k \ge 469.4
\]

\[
k \ge 470
\]

Since:

\[
k \sim \text{Binomial}(1000,0.52)
\]

\[
P(k\ge470)
=
\sum_{k=470}^{1000}
\binom{1000}{k}
(0.52)^k
(0.48)^{1000-k}
\]

\hrule
\vspace{0.5cm}

\section*{Q8}

\[
X \sim \text{Exponential}\left(\frac{1}{2}\right)
\]

PDF:

\[
f(x)=\frac{1}{2}e^{-x/2}
\]

\subsection*{(a)}

\[
P(X>2)=e^{-1}\approx0.3679
\]

\subsection*{(b)}

\[
P(X\ge10|X>9)
=
\frac{e^{-5}}{e^{-9/2}}
=
e^{-1/2}
\]

\hrule
\vspace{0.5cm}

\section*{Q9}

Define:

\[
Z=\sqrt{X^2+Y^2}
\]

\[
F_Z(z)=P(X^2+Y^2\le z^2)
\]

Region:

\[
x^2+y^2\le z^2
\]

CDF:

\[
F_Z(z)
=
\int_{-z}^z
\int_{-\sqrt{z^2-x^2}}^{\sqrt{z^2-x^2}}
f_X(x)f_Y(y)\,dy\,dx
\]

Using Leibniz rule:

\[
\frac{d}{dz}\sqrt{z^2-x^2}
=
\frac{z}{\sqrt{z^2-x^2}}
\]

Final PDF:

\[
f_Z(z)
=
z
\int_{-z}^z
\frac{
f_X(x)
\left[
f_Y(\sqrt{z^2-x^2})
+
f_Y(-\sqrt{z^2-x^2})
\right]
}{
\sqrt{z^2-x^2}
}
dx,
\quad z>0
\]

\hrule
\vspace{0.5cm}

\section*{Q10}

Define:

\[
M_n=\max(X_1,\dots,X_n)
\]

\[
P(M_n\le x)=F(x)^n
\]

\subsection*{Joint CDF}

\subsubsection*{Case 1: $x \le y$}

\[
P(M_n\le x, M_{n+1}\le y)
=
F(x)^n F(y)
\]

\subsubsection*{Case 2: $x>y$}

\[
P(M_n\le x, M_{n+1}\le y)
=
F(y)^{n+1}
\]

\subsection*{Final Answer}

\[
P(M_n\le x, M_{n+1}\le y)
=
\begin{cases}
F(x)^n F(y), & x\le y \\
F(y)^{n+1}, & x>y
\end{cases}
\]

\end{document}